\documentclass[10pt,twocolumn]{article}

%% Packages
\usepackage[margin=0.75in,top=0.9in,bottom=0.9in]{geometry}
\usepackage{times}
\usepackage{amsmath,amssymb}
\usepackage{graphicx}
\usepackage{booktabs}
\usepackage{hyperref}
\usepackage{microtype}
\usepackage{enumitem}
\usepackage{fancyhdr}
\usepackage[numbers,sort&compress]{natbib}

%% Header / footer
\pagestyle{fancy}
\fancyhf{}
\renewcommand{\headrulewidth}{0pt}
\fancyhead[C]{\small\textit{Preprint -- submitted to arXiv, September 2026}}
\fancyfoot[C]{\thepage}

%% Title
\title{\textbf{Recall-Optimised Failure Detection in Industrial Telemetry:}\\
\textbf{A Rolling Degraded-State Counter Versus a 21-Sensor Random Forest}\\
\textbf{on NASA C-MAPSS FD001}}

\author{Puru Pandey\\
\small Department of Artificial Intelligence \& Machine Learning\\
\small Galgotias University, Greater Noida, UP 201310, India\\
\small \href{mailto:purupandey2001@gmail.com}{purupandey2001@gmail.com}\\
\small \href{https://github.com/Puru2001pandey/industrial-telemetry-analytics-research}{github.com/Puru2001pandey/industrial-telemetry-analytics-research}}

\date{September 2026}

\begin{document}
\maketitle
\thispagestyle{fancy}

%% Abstract
\begin{abstract}
Binary failure classification in predictive maintenance requires detecting at-risk
machine states before failure occurs.
Evaluation metrics that aggregate across classes---such as weighted F1---can mask poor
performance on the minority at-risk class, which is the class that matters for
deployment.
This paper compares two classifiers on the NASA C-MAPSS FD001 benchmark (20,631 cycles,
100 engine units): (1) logistic regression trained on a single rolling degraded-state
counter derived from per-unit-normalised sensor composites, and (2) a 100-tree Random
Forest trained on all 21 raw sensor channels.
Both models are evaluated on a unit-level 80/20 split to prevent temporal leakage.
The logistic regression achieves \textbf{0.98 recall} on at-risk cycles (within 30
cycles of failure) versus \textbf{0.88} for the Random Forest, while the Random Forest
achieves higher weighted F1 (0.9576 vs.\ 0.9247).
In the deployment scenario, where a missed failure causes unplanned downtime and a
false alarm triggers only a preventive inspection, at-risk recall is the appropriate
primary metric.
We further show that the recall advantage of the single-feature model is stable across
window lengths of 5--20 cycles and degradation-percentile thresholds of 60--80\%,
suggesting it is not an artefact of a lucky parameter choice.
Code and data are available at
\url{https://github.com/Puru2001pandey/industrial-telemetry-analytics-research}.
\end{abstract}

\noindent\textbf{Keywords:} predictive maintenance, feature engineering, C-MAPSS,
logistic regression, random forest, at-risk recall, class imbalance, industrial IoT.

%%===========================================================================
\section{Introduction}
\label{sec:intro}

Predictive maintenance (PdM) systems aim to detect incipient failures in industrial
equipment before they occur \cite{jardine2006review}.
The NASA C-MAPSS benchmark \cite{saxena2008damage} has become a standard testbed for
prognostic methods, with most prior work focusing on Remaining Useful Life (RUL)
regression \cite{li2018remaining,wu2021temporal}.
However, real-world maintenance decisions often require a binary alert: will this
machine fail within the next $N$ cycles?

For this binary task, the cost structure is asymmetric.
A missed failure---a false negative---leads to unplanned downtime, safety incidents,
and consequential damage costs that dwarf the cost of a scheduled inspection.
A false alarm---a false positive---triggers an unnecessary inspection and minor
operational disruption.
The standard aggregate evaluation metric, weighted F1, treats these two error types
symmetrically and thus may favour models that are suboptimal for deployment: a model
with high weighted F1 on a class-imbalanced dataset can still miss a large fraction of
the failures it is meant to detect.

This paper examines the extent to which a single, domain-informed temporal feature
can improve failure recall compared to a high-dimensional sensor ensemble on C-MAPSS
FD001.
Specifically, we compare logistic regression trained on a rolling count of
degraded operating cycles against a 100-tree Random Forest trained on all 21 raw
sensor channels.
The primary evaluation criterion is recall on the at-risk class (engines within 30
cycles of failure).
We find that the single-feature model achieves 0.98 recall versus 0.88 for the ensemble
while the ensemble achieves higher weighted F1 (0.9576 vs.\ 0.9247)---and argue that
this result has direct implications for how PdM models should be evaluated and selected
for industrial deployment.


%%===========================================================================
\section{Related Work}
\label{sec:related}

The C-MAPSS benchmark \cite{saxena2008damage} has attracted a large body of work, most
of it aimed at RUL regression.
Li et al.\ \cite{li2018remaining} achieve strong RMSE on FD001 with deep convolutional
networks; Wu et al.\ \cite{wu2021temporal} extend this with attention mechanisms across
sensor channels.
These are impressive results, but they address a different question---predicting how
many cycles remain---rather than the binary question of whether a failure is imminent.

For binary classification specifically, Khelif et al.\ \cite{khelif2017direct} show
that a hand-crafted health indicator can match support vector machines on fixed-horizon
prediction.
Ramasso and Saxena \cite{ramasso2014application} make a point that shaped the
preprocessing decisions in this paper: raw C-MAPSS sensor readings carry substantial
engine-specific baseline offsets, so any meaningful comparison requires per-unit
normalisation before modelling.
Without it, differences between engines swamp the degradation signal.

What I could not find in the literature is a direct, cleanly-controlled comparison
between a single domain-informed feature and a full-sensor ensemble, evaluated
specifically on at-risk recall rather than RMSE, accuracy, or overall F1.
Papers that use simpler features tend to compare against other simple features.
Papers that use ensemble methods tend to optimise aggregate metrics.
The narrow question---does one number derived from expert knowledge of degradation
beat 21 raw channels on the metric that deployment actually requires?---does not seem
to have been asked this directly.
That gap is what this paper fills.



%%===========================================================================
\section{Dataset}
\label{sec:data}

\paragraph{NASA C-MAPSS FD001.}
The Commercial Modular Aero-Propulsion System Simulation (C-MAPSS)
dataset~\cite{saxena2008damage} provides run-to-failure telemetry for turbofan engines
simulated under a single operating condition (FD001 subset).
Each row records one operational cycle for one engine unit: three operating-setting
columns and 21 sensor readings.
The training file contains 20,631 rows across 100 engine units, with per-unit cycle
counts ranging from 128 to 362 (mean 206).
No holdout test file is used; evaluation is performed on a held-out partition of
training engines.

\paragraph{Binary label definition.}
For engine $u$, let $C_u^{\max}$ denote its maximum recorded cycle (the failure point).
The Remaining Useful Life at cycle $c$ is $\mathrm{RUL}(u,c)=C_u^{\max}-c$.
We define:
\begin{equation}
  y_{u,c}=
  \begin{cases}
    1 & \text{if }\mathrm{RUL}(u,c)\leq 30 \quad\text{(at risk)}\\
    0 & \text{otherwise}\quad\text{(healthy)}
  \end{cases}
\end{equation}
This yields 17,531 healthy and 3,100 at-risk observations (14.4\% minority class).

%%===========================================================================
\section{Methodology}
\label{sec:method}

\subsection{Sensor Preprocessing}

Several C-MAPSS sensors are constant or near-constant across FD001 and provide no
discriminative signal \cite{ramasso2014application}.
We retain 14 sensors with meaningful variance:
$\{s_{02},s_{03},s_{04},s_{07},s_{08},s_{09},s_{11},s_{12},s_{13},s_{14},s_{15},s_{17},s_{20},s_{21}\}$,
categorised by degradation direction:
rising sensors $\mathcal{R}=\{s_{02},\ldots,s_{15}\}$ (values increase with degradation)
and falling sensors $\mathcal{F}=\{s_{17},s_{20},s_{21}\}$ (values decrease).

Per-unit min-max normalisation removes engine-specific baseline offsets:
\begin{equation}
  \tilde{x}_{u,c,s}=
  \frac{x_{u,c,s}-\min_c x_{u,c,s}}
       {\max_c x_{u,c,s}-\min_c x_{u,c,s}+\varepsilon}
\end{equation}
where $\varepsilon=10^{-9}$ prevents division by zero.

\subsection{Degradation Score and Health-Window Feature}

We define a composite degradation score for each cycle:
\begin{equation}
  d_{u,c}=
  \frac{
    \tfrac{1}{|\mathcal{R}|}\!\sum_{s\in\mathcal{R}}\tilde{x}_{u,c,s}
    +\!\left(1-\tfrac{1}{|\mathcal{F}|}\!\sum_{s\in\mathcal{F}}\tilde{x}_{u,c,s}\right)
  }{2}
\end{equation}

A cycle is labelled \emph{degraded} if its score exceeds the 70th percentile of that
engine's score distribution:
\begin{equation}
  g_{u,c}=\mathbf{1}\!\left[d_{u,c}\geq Q_{0.70}(\{d_{u,c'}\}_{c'})\right]
\end{equation}

The \textbf{health-window feature} is the count of degraded cycles in the 10-cycle
rolling window ending at cycle $c$:
\begin{equation}
  h_{u,c}=\sum_{k=\max(1,\,c-9)}^{c}g_{u,k}
\end{equation}
This scalar encodes the recent degradation trajectory as a single interpretable value.

\subsection{Model A: Logistic Regression}

We train $\ell_2$-regularised logistic regression on the scalar $h_{u,c}$ with
class weights inversely proportional to class frequency (\texttt{balanced}).
This 1-feature model has a trivially interpretable decision rule.

\subsection{Model B: Random Forest Baseline}

A 100-tree Random Forest is trained on all 21 standardised raw sensor columns
(\texttt{class\_weight=`balanced'}, seed 42).

\subsection{Evaluation Protocol}

\paragraph{Unit-level split.}
80 engine units are used for training and 20 for testing (seed 42 shuffle).
All cycles from a given engine appear in exactly one partition, eliminating
temporal leakage from cycle-level random splits.

\paragraph{Primary metric.}
Recall on the at-risk class: the fraction of at-risk cycles correctly detected.
Weighted F1 is reported as a secondary reference.

%%===========================================================================
\section{Results}
\label{sec:results}

Table~\ref{tab:results} summarises performance on 4,291 test cycles from 20 held-out
engine units.

\begin{table}[h]
\centering
\caption{Classification on NASA C-MAPSS FD001 (20 held-out engines, 4,291 test cycles).
At-risk support: 620 cycles. \textbf{Bold}: primary metric.}
\label{tab:results}
\vspace{4pt}
\begin{tabular}{lcc}
\toprule
Model & At-Risk Recall & Weighted F1 \\
\midrule
LR -- health-window (1 feat.) & \textbf{0.98} & 0.9247 \\
RF -- 21 raw sensors           & 0.88          & \textbf{0.9576} \\
\bottomrule
\end{tabular}
\end{table}

\paragraph{Full classification reports.}

\noindent\textit{Logistic Regression (Model A):}
\begin{center}\small
\begin{tabular}{lcccc}
\toprule
Class & Prec. & Recall & F1 & Support\\
\midrule
Healthy  & 1.00 & 0.91 & 0.95 & 3,671\\
At-Risk  & 0.64 & \textbf{0.98} & 0.78 & 620\\
\midrule
Weighted & 0.94 & 0.92 & 0.92 & 4,291\\
\bottomrule
\end{tabular}
\end{center}

\noindent\textit{Random Forest (Model B):}
\begin{center}\small
\begin{tabular}{lcccc}
\toprule
Class & Prec. & Recall & F1 & Support\\
\midrule
Healthy  & 0.98 & 0.97 & 0.97 & 3,671\\
At-Risk  & 0.83 & 0.88 & 0.86 & 620\\
\midrule
Weighted & 0.96 & 0.96 & 0.96 & 4,291\\
\bottomrule
\end{tabular}
\end{center}

\paragraph{Confusion matrices.}

\begin{table}[h]
\centering
\caption{Confusion matrices on 4,291 test cycles.
LR misses 12 at-risk cycles; RF misses 74.}
\label{tab:cm}
\vspace{4pt}
\small
\begin{tabular}{lcc|cc}
\toprule
 & \multicolumn{2}{c|}{LR} & \multicolumn{2}{c}{RF} \\
 & Pred.\ 0 & Pred.\ 1 & Pred.\ 0 & Pred.\ 1 \\
\midrule
Actual 0 & 3{,}329 & 342 & 3{,}559 & 112 \\
Actual 1 & 12 & 608 & 74 & 546 \\
\bottomrule
\end{tabular}
\end{table}

%%===========================================================================
\subsection{Sensitivity Analysis}
\label{sec:sensitivity}

The window length $w$ and the degradation percentile threshold $\tau$ are design
choices.
To assess whether the recall advantage of the single-feature model depends on a
specific parameter setting, we vary $w \in \{5,10,15,20\}$ and
$\tau \in \{50,60,70,80,90\}$.
For within-analysis consistency, both models receive per-unit min-max normalised
inputs in this section (the RF therefore receives normalised rather than globally
standardised sensors, causing its recall to differ slightly from
Table~\ref{tab:results}).

\begin{table}[h]
\centering
\caption{Sensitivity to window $w$ ($\tau=70\%$). LR recall exceeds RF recall at
all settings.}
\label{tab:sens_window}
\vspace{4pt}
\small
\begin{tabular}{ccccc}
\toprule
$w$ & LR Rec. & LR Prec. & RF Rec. & RF Prec. \\
\midrule
 5 & 0.947 & 0.625 & 0.937 & 0.815 \\
10 & 0.976 & 0.643 & 0.937 & 0.815 \\
15 & 0.982 & 0.677 & 0.937 & 0.815 \\
20 & 0.989 & 0.668 & 0.937 & 0.815 \\
\bottomrule
\end{tabular}
\end{table}

\begin{table}[h]
\centering
\caption{Sensitivity to percentile threshold $\tau$ ($w=10$). Higher $\tau$ reduces
flagged cycles, increasing precision and slightly reducing recall.
$\tau=80\%$ approaches RF precision while maintaining recall advantage.}
\label{tab:sens_pct}
\vspace{4pt}
\small
\begin{tabular}{ccccc}
\toprule
$\tau$ & LR Rec. & LR Prec. & RF Rec. & RF Prec. \\
\midrule
50\% & 0.989 & 0.540 & 0.937 & 0.815 \\
60\% & 0.986 & 0.573 & 0.937 & 0.815 \\
70\% & 0.976 & 0.643 & 0.937 & 0.815 \\
80\% & 0.952 & 0.798 & 0.937 & 0.815 \\
90\% & 0.960 & 0.745 & 0.937 & 0.815 \\
\bottomrule
\end{tabular}
\end{table}

Across all 20 parameter combinations, LR recall (0.947--0.989) exceeds RF recall
(0.937).
The recall advantage is not an artefact of a single lucky parameter setting.
The $\tau=80\%$ setting provides a practically useful operating point for deployments
with higher inspection costs: at-risk precision rises to 0.798 (close to RF's 0.815)
while recall remains 1.5 points above the RF (0.952 vs.\ 0.937).

%%===========================================================================
\section{Discussion}
\label{sec:discussion}

The central finding is the confusion matrix.
The Random Forest misses 74 at-risk cycles out of 620; the logistic regression misses
12.
Both models were evaluated on the same held-out engines with balanced class weights.
The Random Forest has higher precision on at-risk predictions (0.83 vs.\ 0.64) and
higher weighted F1 overall---by these aggregate measures it looks like the better
model.
But 74 undetected imminent failures versus 12 is not a small difference in a
safety-critical setting.
Translating to a rough cost model: if a missed failure costs \$50,000 (unplanned
downtime) and a false alarm costs \$5,000 (unnecessary inspection), the Random Forest
incurs an expected failure cost of $74 \times \$50{,}000 = \$3{,}700{,}000$ from missed
detections alone, against $12 \times \$50{,}000 = \$600{,}000$ for the logistic
regression.
The logistic regression generates more false alarms (342 vs.\ 112), adding $342 \times
\$5{,}000 = \$1{,}710{,}000$ versus $112 \times \$5{,}000 = \$560{,}000$.
Net expected cost: LR \$2.31M vs.\ RF \$4.26M---the simpler model is preferable under
any cost ratio where failures cost substantially more than inspections.

The interpretability argument is secondary but relevant.
The health-window feature $h_{u,c} \in \{0,\ldots,w\}$ is a count of recently degraded
cycles.
The decision rule reduces to a threshold on a single integer, communicable to
maintenance staff without ML expertise.
The Random Forest's decision process across 100 trees and 21 input channels is not
interpretable in the same sense.
In industrial settings where operators need to understand and trust an alert before
acting on it, interpretability affects adoption \cite{jardine2006review}.

The precision gap (0.64 for LR vs.\ 0.83 for RF) is a genuine cost.
The sensitivity analysis suggests it is partially tunable: at $\tau=80\%$, LR
precision rises to 0.798 while recall remains above the RF baseline (0.952 vs.\
0.937).
Practitioners should select $\tau$ based on the inspection-to-failure cost ratio of
their specific deployment.

\subsection{Limitations}

This study has several limitations that constrain the generality of its conclusions.
First, FD001 is the simplest C-MAPSS subset---single operating condition, single fault
mode, 100 engines with relatively uniform degradation trajectories.
The per-unit normalisation and degradation composite that work well here may not
transfer cleanly to FD002--FD004, which have multiple operating regimes requiring
explicit condition clustering before health indicators can be computed.
Second, neither the window length $w$ nor the percentile threshold $\tau$ is derived
from physical principles; both are design choices validated empirically within this
dataset.
A learned threshold---derived from, say, a one-class SVM or an isolation forest on
healthy-only training data---would be more principled.
Third, the 30-cycle at-risk horizon is a modelling assumption; different applications
may require shorter or longer detection horizons, and recall at horizon 30 may not
translate to recall at horizon 15 or 50.
Finally, no comparison is made against other simple baselines (single-sensor
thresholds, moving average of raw composite score), which would clarify whether the
rolling-count feature specifically, or simply any low-dimensional degradation
summary, drives the recall advantage.





%%===========================================================================
\section{Conclusion}
\label{sec:conclusion}

A single rolling degraded-state counter gives a logistic regression 0.98 recall on
at-risk engine cycles in NASA C-MAPSS FD001.
A Random Forest on all 21 sensor channels gets 0.88.
The Random Forest has higher weighted F1.

I think the right way to read this result is not ``domain features always win'' but
``evaluation metric choice is a design decision with operational consequences.''
If you deploy a model that maximises weighted F1 on a 14\% minority-class dataset,
you are likely deploying a model that performs well on paper and misses failures in
practice.
The choice of what to optimise should come from the cost structure of the problem,
not from which number looks best in a table.

The full pipeline is available on GitHub under the MIT licence.
Clone the repository, run the Phase~3B notebook cells, and the numbers in this paper
reproduce exactly.



%%===========================================================================
\section*{Code Availability}
All code, preprocessing pipeline, and the NASA C-MAPSS FD001 training file are
available at:
\url{https://github.com/Puru2001pandey/industrial-telemetry-analytics-research}
Results are reproduced by running the Phase~3B notebook cells.

\section*{Acknowledgements}
The author thanks the NASA Ames Prognostics Center of Excellence for making the
C-MAPSS dataset publicly available.
This work was conducted independently, without institutional funding.

%%===========================================================================
\bibliographystyle{plainnat}
\begin{thebibliography}{8}

\bibitem{jardine2006review}
Jardine, A.K.S., Lin, D., and Banjevic, D. (2006).
A review on machinery diagnostics and prognostics implementing condition-based maintenance.
\textit{Mechanical Systems and Signal Processing}, 20(7):1483--1510.

\bibitem{saxena2008damage}
Saxena, A., Goebel, K., Simon, D., and Eklund, N. (2008).
Damage propagation modeling for aircraft engine run-to-failure simulation.
In \textit{Proc.\ Int.\ Conf.\ Prognostics and Health Management (PHM)}, Denver, CO.

\bibitem{ramasso2014application}
Ramasso, E. and Saxena, A. (2014).
Performance benchmarking and analysis of prognostic methods for CMAPSS datasets.
\textit{Int.\ J.\ Prognostics and Health Management}, 5(2):1--15.

\bibitem{li2018remaining}
Li, X., Ding, Q., and Sun, J.-Q. (2018).
Remaining useful life estimation in prognostics using deep convolution neural networks.
\textit{Reliability Engineering \& System Safety}, 172:1--11.

\bibitem{wu2021temporal}
Wu, J. et al.\ (2021).
Data-driven remaining useful life prediction via multiple sensors with attention mechanism.
\textit{Reliability Engineering \& System Safety}, 208:107340.

\bibitem{heimes2008recurrent}
Heimes, F.O. (2008).
Recurrent neural networks for remaining useful life estimation.
In \textit{Proc.\ Int.\ Conf.\ Prognostics and Health Management (PHM)}, Denver, CO.

\bibitem{khelif2017direct}
Khelif, R. et al.\ (2017).
Direct remaining useful life estimation based on support vector regression.
\textit{IEEE Trans.\ Industrial Electronics}, 64(3):2276--2285.

\bibitem{mobley2002maintenance}
Mobley, R.K. (2002).
\textit{An Introduction to Predictive Maintenance}, 2nd ed.
Butterworth-Heinemann.

\end{thebibliography}

\end{document}
